\documentclass[10pt]{article} 
\renewcommand{\thempfootnote}{\arabic{mpfootnote}}

\usepackage[utf8]{inputenc}

\usepackage{multicol}
\setlength{\columnsep}{1cm}

\usepackage{abc}
\usepackage{graphicx}
\usepackage{musicography}

\usepackage[margin=0.5in]{geometry} 
\geometry{letterpaper}

\usepackage{fancyhdr} 
\pagestyle{fancy} % options: empty , plain , fancy
\renewcommand{\headrulewidth}{0pt} % customise the layout...
\lhead{}\chead{}\rhead{}
\lfoot{}\cfoot{\thepage}\rfoot{}

\begin{document}

\thispagestyle{empty}

\begin{abc}[name=frostymorning,options={-O= +c -F single.fmt}]
X:1
T:Cold Frosty Morning
M:C|
K:Ador
|:"Am"EDEG A2AA | cAcd e2 A2 | "G"B2 G2 G2-GD | GABc "Em"d2 d2 | 
"Am"EDEG A2AA | cAcd e2 ed | cBAc "G"BAGE | "Am"A2 AA A4 :| 
|: "Am"eg a2 a2 aa | abag e2 e2 | "Em"edef gagf | "G"edcd "Em"e2 e2 | 
"Am"A2 AA "C"c2 cc | "D"dcd2 "E"e2 ed | "Am"cBAc "G"BAGE | "Am"A4 A4 :| 
\end{abc}
\vskip 0.5in
\noindent\makebox[\textwidth][c]{%
  \begin{minipage}{7.0in}
   \setlength{\parindent}{1.5em}

   This Old Time tune staple was collected by the folklorist Alan
   Jabbour, from the Virginia fiddler Henry Reed, for the Library of
   Congress.\footnote{http://memory.loc.gov/diglib/ihas/loc.afc.afcreed.13037a09/default.html}
   It is most closely associated with the Appalachian Mountain region,
   Virginia, West Virginia, and Tennessee in particular. Some sources
   state that this tune is originally Scottish as it is said to
   commemorate the Battle of Culloden when on April 16th, 1746, the
   Scots were defeated by the English thus ending the Jacobite
   Rebellion.\footnote{http://communityguitar.com/students/Songs/cold\_frosty\_morning.htm}
   Other sources, including Jabbour, point out that known older
   Scottish tunes have no relation to this one except in the title,
   and that this tune can only be traced to Henry Reed with any
   certainty.\footnote{https://banjonews.com/2010-05/licks\_and\_letters.html}
   Did Reed learn it from an older tradition, or is it original
   to him? Perhaps we shall never know, friends.
   
 \end{minipage}
}


\end{document}
